\addcontentsline{toc}{section}{Заключение}

В ходе проведенного исследования был осуществлен сравнительный аналази современных технологий разработки интеллектуальных систем: фреймовых языков формализации данных и корпоративных платформ для развертывания интеллектуальных систем.

В первом разделе работы были подробно рассмотрены фреймовые языки FrameNet и KL-ONE. Проанализированы их синтаксические и семантические особенности,  выявлены общие черты и различия в подходах к представлению и обработке данных. Результаты анализа были формализованы с помощью SCn-кода.

Во втором разделе были подробно рассмотрены инструменты для разработки мультиагентных систем. Проведен анализ их функциональных возможностей, методов коммуникации агентов,их создания, а также особенностей каждого инструмента. Информация о каждом инструменте также была формализована на SCn-коде.

Результаты работы могут быть использованы при проектировании корпоративных интеллектуальных систем, для которых используется фреймовый подход к формализации данных, многоагентная структура, а также при выборе и обосновании конкретных технологических решений для прикладных задач.